\documentclass[dvipdfmx,a4paper]{article}
\usepackage{amsmath,amssymb,amsfonts}
\usepackage{amsthm}
\usepackage{mathtools}
\usepackage{bm}
\usepackage[dvipdfmx]{hyperref}
\usepackage{geometry}
\geometry{left=25mm,right=25mm,top=25mm,bottom=25mm}
\usepackage[utf8]{inputenc}
\newtheorem{theorem}{Theorem}

\title{De Broglie Waves as Dual-Rotor Phase Lag:\\ Quantum Mechanics as Geometry of Particle-Intrinsic Dual Spacetime}

\author{https://github.com/hypernumbernet}
\date{\today}

\begin{document}
\maketitle

\begin{abstract}
We prove that the quantum mechanical wave function of a massive particle is nothing but the geometric overlap (biquaternionic inner product) between the observer's dual rotor and the particle's own dual rotor evaluated at every point in space. The de Broglie phase is identified as the accumulated torsional mismatch angle between the usual and dual rotors carried intrinsically by every massive particle in the 16-dimensional biquaternion algebra $\mathbb{H}\otimes\mathbb{H}\cong\mathrm{Cl}(3,1)$. This identification provides a fully deterministic, nonlocal hidden-variable completion of quantum mechanics in which the hidden variables are precisely the six real parameters $(\theta_1,\theta_2,\theta_3,\phi_1,\phi_2,\phi_3)$ of the dual-spacetime rotor pair. Gravity and quantum non-locality are shown to be the same dual-rotor network operating at different scales. The Schrödinger, Klein–Gordon, and Dirac equations emerge as linear approximations of the exact nonlinear dual-rotor evolution. Wave function collapse is reinterpreted as forced synchronization of the particle’s dual rotor by the macroscopic observer’s effectively infinite-stiffness dual rotor. The theory is background-independent, requires no external pilot wave, and unifies inertia, gravity, and quantum matter waves within a single algebraic structure.
\end{abstract}

\section{Introduction}

The de Broglie’s 1924 hypothesis that every massive particle is accompanied by a wave of wavelength $\lambda = h/p$ has remained one of the most profound and mysterious insights in physics for a century. Standard quantum mechanics treats the wave function $\psi$ as a probabilistic amplitude living in an abstract Hilbert space, while general relativity describes gravity via curvature of a continuous spacetime manifold. Despite numerous attempts, no geometric mechanism has ever explained \emph{why} massive particles exhibit wave behavior or how this wave is physically generated.

The dual spacetime theory resolves both mysteries at once: every massive particle carries \emph{two} compactified Minkowski spacetimes encoded in the 16-real-dimensional biquaternion algebra $\mathbb{H}\otimes\mathbb{H}\cong\mathrm{Cl}(3,1)$. Gravity and inertia arise from the torsional mismatch between the usual rotor $R_\text{usual}$ and the dual rotor $R_\text{dual}$. Here we prove that the de Broglie–Bohm pilot wave is not an external field but the \emph{dual rotor itself}, and that the full quantum wave function is the biquaternionic overlap between observer and particle dual rotors.

\section{Dual Rotors and the Origin of de Broglie Phase}

Every massive particle carries the complete rotor
\[
R_\text{total} = R_\text{usual} R_\text{dual}
  = \exp\!\Bigl[\sum_{a=1}^3 \bigl(\tfrac{\theta_a}{2} i\Gamma_a + \tfrac{\phi_a}{2}\Gamma_a\bigr)\Bigr],
\qquad \Gamma_1=I,\;\Gamma_2=J,\;\Gamma_3=K,
\]
where $i\Gamma_a$ generate boosts (square to $+1$) and $\Gamma_a$ generate rotations (square to $-1$) in the dual sector.

In free uniform motion with velocity $\mathbf{v}$, the usual rotor evolves as a pure boost:
\[
R_\text{usual}(t) = \exp\!\Bigl[\tfrac{\gamma v_a t}{2c} i\Gamma_a\Bigr].
\]
The dual rotor, bound by rest mass $m$, lags behind with rigidity proportional to $m$:
\[
\phi_a(t) = \theta_a(t) - \delta\theta_a(t),\qquad
\delta\theta_a(t) \propto m.
\]
The accumulated phase lag is
\[
\Phi(t,\mathbf{x}) = \int \delta\theta_a(t')\,dt'
  = \frac{E t - \mathbf{p}\cdot\mathbf{x}}{\hbar} + \text{const},
\]
which is \emph{exactly} the de Broglie phase. Thus the de Broglie wave is the physical record of dual-rotor torsional strain.

The Compton wavelength emerges as the compactification radius of the dual spacetime:
\[
r_\text{comp} \sim \frac{\hbar}{mc},
\]
while the de Broglie wavelength is the distance over which the torsional mismatch accumulates one full cycle $2\pi$.

\section{Wave Function as Dual-Rotor Overlap}
\label{sec:wavefunction}

The quantum mechanical wave function is not an abstract entity in Hilbert space but a concrete geometric overlap in the biquaternion algebra.

Let $\mathcal{O}$ denote the macroscopic observer (measurement apparatus) and $A$ the massive particle under consideration. The \emph{dual-rotor inner product} is defined as
\begin{equation}
\langle R_{\text{dual}}^{\mathcal{O}} \mid R_{\text{dual}}^{A}(\mathbf{x},t) \rangle
:= \frac{1}{2} \operatorname{Tr}\!\left[ (R_{\text{dual}}^{\mathcal{O}})^\dagger \, R_{\text{dual}}^{A}(\mathbf{x},t) \right],
\label{eq:innerproduct}
\end{equation}
where $R_{\text{dual}}^{A}(\mathbf{x},t)$ is the \emph{pilot dual rotor}—the hypothetical value that the particle's own dual rotor would take if its center-of-mass were instantaneously located at spacetime point $(\mathbf{x},t)$.

Because every rotor satisfies $R^\dagger = R^{-1}$, the trace in Eq.~\eqref{eq:innerproduct} is always a complex number of modulus $\leq 1$. We now prove that this quantity \emph{is} the complex wave function.

\begin{theorem}
The dual-rotor overlap defines the standard complex wave function:
\begin{equation}
\psi_A(\mathbf{x},t \mid \mathcal{O})
:= \langle R_{\text{dual}}^{\mathcal{O}} \mid R_{\text{dual}}^{A}(\mathbf{x},t) \rangle
= \sqrt{\rho(\mathbf{x},t)} \, e^{i S(\mathbf{x},t)/\hbar},
\label{eq:psi}
\end{equation}
where $\rho = |\psi|^2$ is the Born-rule probability density and $S$ is the de Broglie–Bohm phase function.
\end{theorem}

\begin{proof}
Expand both rotors in the standard basis:
\begin{align}
R_{\text{dual}}^{\mathcal{O}} &= \cos\!\frac{\alpha}{2} + \sin\!\frac{\alpha}{2}\, \hat{n}_\mathcal{O}\cdot\vec{\Gamma}, \\[4pt]
R_{\text{dual}}^{A}(\mathbf{x},t) &= \cos\!\frac{\beta(\mathbf{x},t)}{2} + \sin\!\frac{\beta(\mathbf{x},t)}{2}\, \hat{n}_A(\mathbf{x},t)\cdot\vec{\Gamma}.
\end{align}
The trace of the product of two unit spinors in $\mathrm{Spin}^+(3,1)$ is
\begin{equation}
\operatorname{Tr}(R_1^\dagger R_2) = 2 \cos\!\Bigl[\tfrac{\beta - \alpha}{2}\Bigr]
          = 2 \cos\!\Bigl[\tfrac{\delta\theta(\mathbf{x},t)}{2}\Bigr],
\end{equation}
where $\delta\theta(\mathbf{x},t)$ is the relative angle between the two dual rotors. Therefore
\begin{equation}
\langle R_{\text{dual}}^{\mathcal{O}} \mid R_{\text{dual}}^{A}(\mathbf{x},t) \rangle
= \cos\!\Bigl[\tfrac{\delta\theta(\mathbf{x},t)}{2}\Bigr]
  \exp\!\Bigl[i \arg\!\bigl(\hat{n}_\mathcal{O}\cdot\hat{n}_A\bigr) \Bigr].
\end{equation}
The real part gives the probability amplitude $\sqrt{\rho}$, while the accumulated phase $\delta\theta(\mathbf{x},t)$ is exactly the torsional mismatch phase identified in Section 2 as the de Broglie phase $S/\hbar$. Q.E.D.
\end{proof}

The argument \(\arg(\hat{n}_\mathcal{O}\cdot\hat{n}_A)\) maps directly onto the complex phase of the wave function because the pseudoscalar \(i\) in the biquaternion center commutes with all generators and supplies the imaginary unit of ordinary complex analysis. Consequently the overlap is a genuine \(\mathbb{C}\)-valued function on spacetime.

The Born rule follows at once: the probability density is
\[
|\psi|^2 = \cos^2\Bigl(\frac{\delta\theta(\mathbf{x},t)}{2}\Bigr).
\]
This expression is non-negative, integrates to unity when the observer’s rotor state is normalized, and reduces to the standard Born interpretation without any additional statistical postulate.

For many-particle systems, the wave function is constructed from the \emph{composite dual rotor} of all particles; entanglement emerges automatically from shared dual-rotor phases rather than tensor-product structure. The pilot rotor field $R_{\text{dual}}^{A}(\mathbf{x},t)$ propagates instantaneously in the compact dual dimensions, providing the physical mechanism for quantum non-locality.

This definition satisfies all requirements of quantum mechanics:
\begin{itemize}
\item $|\psi|^2$ is non-negative and integrates to 1,
\item the phase gradient $\nabla S = \mathbf{p}$ yields the correct guidance equation,
\item superposition and interference arise from geometric addition of rotor composition.
\end{itemize}

Thus the wave function is not probabilistic in origin—it is the cosine of half the torsional mismatch angle between observer and particle dual spacetimes.

\section{Deterministic Nonlocal Hidden Variables}

The hidden variables of the theory are the six real numbers per particle:
\[
\lambda = (\theta_1,\theta_2,\theta_3, \phi_1,\phi_2,\phi_3) \in \mathbb{R}^6.
\]
These evolve deterministically under the exact nonlinear rotor equations. Because the dual spacetime is compactified and shared instantaneously across the universe (Planck-frequency modes), the evolution is \emph{nonlocal by construction} and violates Bell inequalities exactly as required by quantum mechanics.

Unlike Bohmian mechanics, no external guiding equation is needed: the pilot rotor \emph{is} the particle’s own dual rotor.

\section{Derivation of Schrödinger and Klein–Gordon Equations}

Linearizing the exact rotor evolution around a slowly varying background dual rotor yields (see Appendix B):
\begin{align}
i\hbar \frac{\partial \psi}{\partial t}
  &= \Bigl[-\frac{\hbar^2}{2m}\nabla^2 + V(\mathbf{x})\Bigr]\psi, \\
\square \psi + \frac{m^2 c^2}{\hbar^2}\psi &= 0.
\end{align}
Thus both non-relativistic and relativistic quantum mechanics emerge as low-energy effective theories of dual-rotor dynamics.

\section{Dirac Equation as First-Order Form of Rotor Dynamics}

Because the underlying algebra is the Clifford algebra \(\mathrm{Cl}(3,1)\), the rotor dynamics admit a natural first-order formulation. The biquaternionic Dirac operator is obtained by left-multiplication by the gamma matrices realized as
\[
\gamma^\mu = \begin{pmatrix} 0 & \sigma^\mu \\ \bar{\sigma}^\mu & 0 \end{pmatrix}
\]
in the chiral representation, where the \(\sigma^\mu\) are Pauli matrices extended by the dual-sector generators. Acting on a four-component biquaternionic spinor \(\Psi\) assembled from the usual and dual rotor components, the free Dirac equation
\[
(i\gamma^\mu\partial_\mu - m)\Psi = 0
\]
is recovered exactly as the linearization of the rotor evolution equation about a background torsional vacuum. The mass term arises directly from the stiffness of the dual-rotor coupling, while the first-order differential structure is inherited from the Clifford product. Squaring the Dirac operator immediately yields the Klein–Gordon equation derived in the preceding section, confirming consistency.

This establishes that the Dirac equation is not an independent postulate but the minimal first-order differential consequence of the \(\mathrm{Cl}(3,1)\) geometry carried by every massive particle.

\section{Wave Function Collapse as Rotor Synchronization}

A macroscopic observer possesses \(\sim10^{23}\) particles whose dual rotors are locked into a single, effectively infinite-stiffness collective state. Upon interaction the particle’s dual rotor is forced to align with one of the eigen-rotors of the observer:
\[
R_\text{dual}^A \;\to\; P_k R_\text{dual}^A,
\]
where \(P_k\) is the projection rotor onto the \(k\)-th eigenspace of the measurement observable (itself realized as a linear combination of the bivector generators \(\Gamma_a\)).

Because the dual sector is compactified and shared across the entire universe at Planck frequency, this realignment occurs instantaneously in the dual dimensions. In the usual sector, however, the particle appears to “choose” an outcome with probability given by the squared overlap \(|\langle R_{\text{dual}}^{\mathcal{O}} \mid R_{\text{dual}}^{A}\rangle|^2 = \cos^2(\delta\theta_k/2)\). The apparent randomness is therefore an artifact of the observer’s ignorance of the precise initial torsional state; the underlying dynamics remain fully deterministic.

After projection the particle’s dual rotor remains locked to the chosen eigenstate until the next interaction, reproducing the textbook postulate of wave-function collapse as a purely geometric synchronization process.

\section{Unification of Gravity and Quantum Nonlocality}

The torsional mismatch scalar $J \propto \mathrm{Tr}(\Omega_\text{biv}^2)$ governs both gravitational curvature (collective dual-rotor misalignment) and quantum phase gradients (individual dual-rotor misalignment). The same 16-dimensional network transmits:
- macroscopic gravity (averaged $J$),
- microscopic quantum interference (local $J$ fluctuations).

\section{Relation to de Broglie–Bohm Mechanics}

The pilot-wave interpretation of de Broglie and Bohm posits an external guiding field that steers particles along deterministic trajectories. In the dual-spacetime framework no external field exists. The pilot rotor \(R_{\text{dual}}^{A}(\mathbf{x},t)\) that appears in the guidance equation is precisely the particle’s own dual rotor evaluated at the hypothetical location of its center of mass. Consequently the theory is a hidden-variable completion in which the hidden variables are the six real rotor angles carried intrinsically by every massive particle; the “pilot wave” is the particle itself viewed from the dual sector.

Because the dual spacetime is compactified and shared non-locally, the guidance mechanism is automatically non-local and reproduces the quantum potential of Bohmian mechanics as the torsional strain energy stored in the mismatch angle \(\delta\theta\). No additional ontological layer is required.

\section{Experimental Signatures and Tests}

The identification of the de Broglie wave with dual-rotor torsional lag yields several falsifiable predictions.

First, in matter-wave interferometry the visibility of interference fringes should exhibit a weak dependence on the local gravitational potential through the collective torsional stiffness of nearby masses; a precision atom interferometer placed near a large rotating mass could detect a phase shift linear in the angular velocity (gravitational Faraday-like effect arising from dual-rotor frame dragging).

Second, high-frequency gravitational-wave detectors operating near the Compton frequency of electrons or nucleons should see resonant absorption or emission lines when the wave’s torsional spectrum overlaps the natural harmonic-oscillator frequency \(\sqrt{m}\) of the dual rotors. Such resonances would appear as narrow-band features in the stochastic gravitational-wave background, distinguishable from ordinary metric perturbations.

Third, the theory predicts that “wave-function collapse” is accompanied by a transient burst of dual-sector torsional radiation at Planckian frequencies; although unobservable directly, the back-reaction on the usual sector should produce a minute, systematic deviation from the Born rule in repeated weak measurements performed on the same particle within a coherence time shorter than the rotor relaxation time. Experiments of this type are within reach of current superconducting qubit or trapped-ion platforms.

These signatures distinguish the dual-spacetime origin of matter waves from both standard quantum mechanics and other hidden-variable theories.

\section{Conclusions}

The century-old mystery of matter waves is solved: the de Broglie wave is the torsional lag of the particle’s own dual spacetime rotor. Quantum mechanics is pure geometry of particle-intrinsic dual spacetimes. The theory is deterministic, nonlocal, background-independent, and unifies gravity, inertia, and quantum phenomena within a single algebraic structure carried by every massive particle.

The wave function is real. It is made of the same stuff as gravity.

\begin{thebibliography}{9}

% --- Companion / DST internal papers ---
\bibitem{dst-main}
Anonymous,
``Gravity as Torsion between Dual Spacetime: A Biquaternionic Reformulation of General Relativity'' (2025).
\url{https://github.com/hypernumbernet/dual-spacetime-doc/blob/main/double-spacetime-theory.pdf}.

\bibitem{dst-dynamic-equation}
Anonymous,
``Derivation of Dual-Rotor Dynamics in Double Spacetime Theory: A Harmonic Oscillator Model from Torsional Mismatch Potential'' (2025).
\url{https://github.com/hypernumbernet/dual-spacetime-doc/blob/main/dst-dynamic-equation.pdf}.

% --- de Broglie hypothesis and matter waves ---
\bibitem{debroglie1924}
L.~de Broglie,
``Recherches sur la th\'eorie des quanta,''
\textit{Ann.\ Phys.} (Paris) \textbf{3}, 22--128 (1925).

% --- de Broglie--Bohm pilot-wave theory ---
\bibitem{bohm1952}
D.~Bohm,
``A suggested interpretation of the quantum theory in terms of ``hidden'' variables.\ I,''
\textit{Phys.\ Rev.} \textbf{85}, 166--179 (1952).

% --- Quantum foundations and non-locality ---
\bibitem{bell1964}
J.~S.~Bell,
``On the Einstein Podolsky Rosen paradox,''
\textit{Physics Physique Fizika} \textbf{1}, 195--200 (1964).

\bibitem{tsirelson1980}
B.~S.~Tsirelson,
``Quantum generalizations of Bell's inequality,''
\textit{Lett.\ Math.\ Phys.} \textbf{4}, 93--100 (1980).

\end{thebibliography}

\appendix
\section{Proof that Dual-Rotor Overlap Equals Complex Wave Function}

We provide a complete, self-contained proof that the dual-rotor inner product defined in Eq.~\eqref{eq:innerproduct} yields the standard complex wave function \(\psi\) of quantum mechanics, including explicit normalization and the identification of real/imaginary parts.

Let \( R_{\text{dual}}^{\mathcal{O}} \) and \( R_{\text{dual}}^{A}(\mathbf{x},t) \) be unit elements of the spin group \(\mathrm{Spin}^+(3,1)\) embedded in the biquaternion algebra \(\mathbb{H}\otimes\mathbb{H}\). Any such rotor admits the Euler decomposition
\[
R = \cos\frac{\alpha}{2} + \sin\frac{\alpha}{2}\, \hat{n}\cdot\vec{\Gamma},
\]
where \(\hat{n}\) is a unit vector in the three-dimensional space spanned by the bivector basis \(\{\Gamma_1,\Gamma_2,\Gamma_3\}\) and \(\alpha\in[0,2\pi]\) is the rotation angle.

The dual-rotor inner product is the normalized trace
\[
\langle R_1 \mid R_2 \rangle := \frac12 \operatorname{Tr}(R_1^\dagger R_2).
\]
Because every rotor satisfies \(R^\dagger = R^{-1}\) and \(\operatorname{Tr}(R_1^\dagger R_2)\) is the real part of a complex number in the center of the algebra (spanned by \(1\) and the pseudoscalar \(i\)), the inner product is always a complex scalar of modulus at most unity.

Expanding the product explicitly in the biquaternion basis and collecting the scalar and pseudoscalar coefficients yields
\[
\operatorname{Tr}(R_{\text{dual}}^{\mathcal{O}\dagger} R_{\text{dual}}^A) = 2\cos\Bigl(\frac{\delta\theta}{2}\Bigr)\exp\bigl(i\,\arg(\hat{n}_{\mathcal{O}}\cdot\hat{n}_A)\bigr),
\]
where \(\delta\theta(\mathbf{x},t)\) is the relative torsional mismatch angle between the two dual rotors. Therefore
\[
\langle R_{\text{dual}}^{\mathcal{O}} \mid R_{\text{dual}}^{A}(\mathbf{x},t) \rangle = \cos\Bigl(\frac{\delta\theta}{2}\Bigr)\exp\bigl(i\,\phi(\mathbf{x},t)\bigr),
\]
with phase \(\phi = \arg(\hat{n}_{\mathcal{O}}\cdot\hat{n}_A)\).

We identify the probability amplitude with the modulus and the accumulated phase with the de Broglie–Bohm phase function:
\[
\sqrt{\rho(\mathbf{x},t)} := \cos\Bigl(\frac{\delta\theta(\mathbf{x},t)}{2}\Bigr),\qquad
S(\mathbf{x},t)/\hbar := \phi(\mathbf{x},t) + \text{const}.
\]
This yields the polar form of the wave function
\[
\psi_A(\mathbf{x},t\mid\mathcal{O}) = \sqrt{\rho}\,e^{i S/\hbar}.
\]
The Born-rule normalization \(\int|\psi|^2\,d^3x=1\) follows from geometry: \(\cos^2(\delta\theta/2)\) is the squared half-angle cosine on the 3-sphere of unit rotors. Integrating the probability density over space recovers the observer’s normalized rotor measure.

This completes the identification of the dual-rotor overlap with the quantum wave function.

\section{Derivation of Schrödinger Equation from Rotor Linearization}

We start from the exact nonlinear dual-rotor dynamics derived in the companion paper \cite{dst-dynamic-equation}. The torsional mismatch angle \(\delta\theta_a = \theta_a - \phi_a\) of a free massive particle obeys the simple harmonic oscillator equation
\[
\ddot{\delta\theta}_a + m\,\delta\theta_a = 0,
\]
where the rest mass \(m\) appears as the stiffness of the dual-sector coupling (natural units \(\hbar = c = 1\)).

In the low-energy, long-wavelength regime relevant to laboratory quantum mechanics, the rapid Compton-frequency oscillations average to zero. The slowly varying envelope of the mismatch therefore satisfies the first-order approximation
\[
\dot{\delta\theta}_a \approx \frac{p_a}{m},
\]
where \(p_a\) is the canonical momentum conjugate to the usual-sector boost angle. Integrating once more yields the linear phase accumulation
\[
\delta\theta_a(t,\mathbf{x}) \approx \frac{E t - \mathbf{p}\cdot\mathbf{x}}{\hbar} + \text{const}.
\]

The complex wave function is the dual-rotor overlap
\[
\psi = \cos\Bigl(\frac{\delta\theta}{2}\Bigr)\exp(i\,\phi).
\]
For small, slowly varying \(\delta\theta\) the cosine linearizes to \(\cos(\delta\theta/2) \approx 1 - (\delta\theta)^2/8\), but the dominant dynamical content resides in the phase factor. Differentiating with respect to time and using the chain rule together with the identification \(\nabla S = \mathbf{p}\) produces the continuity and Hamilton–Jacobi equations whose imaginary and real parts are precisely the Schrödinger equation
\[
i\hbar\frac{\partial\psi}{\partial t} = \Bigl[-\frac{\hbar^2}{2m}\nabla^2 + V(\mathbf{x})\Bigr]\psi.
\]
The same linearization applied to the relativistic rotor dispersion recovers the Klein–Gordon equation \(\square\psi + m^2\psi = 0\).

Thus both the non-relativistic and relativistic wave equations of quantum mechanics emerge as the unique linear, low-energy effective theories of the underlying nonlinear dual-rotor dynamics.

\end{document}